\chapter{Gauge Invariance and the Necessity of the Vector Potential}

Before beginning the physics, there are two practical matters.  The next
meeting will be skipped while the lecturer is in Belgium, but the course
will continue until all ten lectures have been given.  The handwritten notes
for the evening have also been prepared at short notice; anyone willing to
scan them for the class website will rescue them from a remarkably slow
scanner.

The previous lecture was mathematically dense.  Rather than pretending
otherwise, we shall go through its central calculation again.  The aim of
the course is a genuine theoretical minimum: enough structure to continue
to the next level.  If the mathematics could be removed without losing the
physics, there would be every reason to remove it.  But exact laws require
quantitative statements, and quantitative statements require mathematics.
The notation may become arcane; that is not evidence that physicists are
trying to hide the subject behind priestly language.

Perhaps a simpler formulation of nature remains to be found.  One can, as
Feynman joked, square every equation of physics, add the squares, and set the
sum equal to zero.  That produces one equation, because every square must
then vanish, but it produces no new understanding.  Simplicity means finding
the organizing idea, not merely compressing the typography.  When an
audience member objects to the suggestion that such an idea must come from
someone young, the point is readily conceded: age is irrelevant to insight.

\lectureoverview{We shall first restore the index notation needed to
recognize Lorentz scalars.  We then formulate four principles for a
relativistic theory, repeat the charged-particle derivation in full, and
finally discover gauge invariance twice: as an endpoint term in the action
and as the invariance of the electromagnetic field tensor.}

\section{The Index Language Revisited}

\subsection{Why some indices are upstairs}

Let
\begin{equation}
x^\mu=(t,x,y,z),
\qquad
\mu=0,1,2,3,
\end{equation}
with metric
\begin{equation}
\eta_{\mu\nu}
=
\operatorname{diag}(-1,1,1,1).
\label{eq:sr7-metric}
\end{equation}
Coordinates and coordinate displacements are assigned upper indices:
\begin{equation}
dx^\mu=(dt,dx,dy,dz).
\end{equation}
Derivatives with respect to those coordinates carry lower indices:
\begin{equation}
\partial_\mu
\equiv
\frac{\partial}{\partial x^\mu}.
\label{eq:sr7-gradient}
\end{equation}
Which family is called upper and which lower is a convention, much as the
names left and right are conventions.  Once chosen, however, the convention
must be followed consistently because the two families transform
differently.

\begin{classroomqa}
\audiencequestion{When should an index be written as a superscript, and when should it be a subscript?}
\lecturerresponse{Coordinates, displacements, and objects that transform like them are assigned upper indices.  Derivatives and objects that transform like derivatives are assigned lower indices.  The choice of which family is called upper is conventional, but the distinction between the two transformation laws is not. \lecturetimestamp{00:05:03}}
\end{classroomqa}

For a scalar function \(s(x)\), the total differential is
\begin{equation}
ds
=
\frac{\partial s}{\partial x^\mu}\,dx^\mu
=
\partial_\mu s\,dx^\mu.
\label{eq:sr7-total-differential}
\end{equation}
Written out, this is
\begin{equation}
ds
=
\frac{\partial s}{\partial t}\,dt
+\frac{\partial s}{\partial x}\,dx
+\frac{\partial s}{\partial y}\,dy
+\frac{\partial s}{\partial z}\,dz.
\end{equation}
The compact expression says exactly the same thing.  It is not a new law;
it is a device for seeing the scalar structure without carrying four terms
through every line.

\begin{figure}[htbp]
\centering
\includegraphics[width=0.90\textwidth]{lecture_07_figure_01.png}
\caption{The scalar differential assembled from a lower-index gradient and an upper-index displacement.}
\lectureframe{00:08:20}
\end{figure}

Einstein's convention is now immediate: an index repeated once upstairs and
once downstairs is summed.  The contraction in
\eqref{eq:sr7-total-differential} is therefore a scalar.  There is a deeper
geometric interpretation of vectors and covectors, but that discussion is
postponed until general relativity.

\subsection{Raising an index}

Every four-vector has upper- and lower-index forms related by the metric:
\begin{equation}
A_\mu=\eta_{\mu\nu}A^\nu.
\end{equation}
With \eqref{eq:sr7-metric},
\begin{equation}
A_0=-A^0,
\qquad
A_m=A^m,
\qquad m=1,2,3.
\label{eq:sr7-lowering}
\end{equation}
Thus raising or lowering an index changes the sign of the time component and
does nothing to the spatial components.  The same rule applies to the
derivative:
\begin{equation}
\begin{aligned}
\partial_\mu
&=
\left(
\frac{\partial}{\partial t},
\frac{\partial}{\partial x},
\frac{\partial}{\partial y},
\frac{\partial}{\partial z}
\right),\\[0.5em]
\partial^\mu
&=
\left(
-\frac{\partial}{\partial t},
\frac{\partial}{\partial x},
\frac{\partial}{\partial y},
\frac{\partial}{\partial z}
\right).
\end{aligned}
\label{eq:sr7-raised-gradient}
\end{equation}

\begin{figure}[htbp]
\centering
\includegraphics[width=0.90\textwidth]{lecture_07_figure_02.png}
\caption{The lower- and upper-index derivative operators; only the time component changes sign.}
\lectureframe{00:12:30}
\end{figure}

\begin{classroomqa}
\audiencequestion{Does a partial derivative always carry a lower index, and can it also be raised?}
\lecturerresponse{The operator \(\partial/\partial x^\mu\) is naturally
\(\partial_\mu\), a covector.  Applying the metric gives
\(\partial^\mu=\eta^{\mu\nu}\partial_\nu\); in flat spacetime this merely
reverses the time component. \lecturetimestamp{00:11:00}}
\end{classroomqa}

The invariant interval is the familiar example:
\begin{equation}
x_\mu x^\mu
=
x^2+y^2+z^2-t^2.
\label{eq:sr7-interval}
\end{equation}
A contraction of one upper and one lower index produces a Lorentz scalar.
This simple recognition rule is the reason for carrying the notation so
carefully.

\section{Four Principles for Constructing a Theory}

Suppose an experiment reveals a new phenomenon and we want to formulate a
theory of it.  Four principles organize the search.  They may eventually
need refinement, but every successful fundamental theory presently fits
their pattern.

\subsection{The action principle}

The first demand is an action.  For particle mechanics,
\begin{equation}
I[x]
=
\int dt\,L(x,\dot x),
\label{eq:sr7-particle-action-general}
\end{equation}
and the physical path is stationary under variations that hold its
endpoints fixed:
\begin{equation}
\delta I=0.
\end{equation}
The resulting equations are
\begin{equation}
\frac{d}{dt}
\left(
\frac{\partial L}{\partial\dot x^m}
\right)
=
\frac{\partial L}{\partial x^m}.
\label{eq:sr7-EL}
\end{equation}

The action principle is not an aesthetic ornament.  It is the framework
that relates symmetries to conservation laws.  Time-translation symmetry of
an action yields energy conservation; spatial translations yield momentum
conservation; related symmetries underlie other conserved quantities.  The
action alone is not the conservation law, but it is the structure through
which the symmetry becomes one.

\subsection{Locality}

The second demand is locality.  If one point of a violin string is disturbed,
the neighboring segment responds first.  Its neighbor responds next, and
the disturbance propagates.  The effect can eventually reach the whole
string, but the direct short-time influence is local.

For a particle, locality is encoded by a Lagrangian evaluated along the
trajectory.  It depends on the position and velocity at the same time, not
on arbitrary pairs of remote times.  For a field \(\phi(x)\),
\begin{equation}
I[\phi]
=
\int d^4x\,
\mathcal L\!\left(\phi,\partial_\mu\phi\right).
\label{eq:sr7-local-field-action}
\end{equation}
The density at one event depends on the field and its first derivatives
there.  It does not contain a direct product of the field here with the
field at an unrelated distant point.

\begin{figure}[htbp]
\centering
\includegraphics[width=0.90\textwidth]{lecture_07_figure_03.png}
\caption{Particle and field actions displayed beneath the locality principle.}
\lectureframe{00:19:30}
\end{figure}

This is why local equations of motion are differential equations.
Derivatives compare infinitesimally neighboring values.  A nonlocal
alternative might relate the acceleration now to a force at a completely
different time, or relate velocities at two remote times to a force at a
third.  Such laws are mathematically imaginable, but locality excludes
them.

\begin{classroomqa}
\audiencequestion{Does quantum entanglement contradict this locality principle?}
\lecturerresponse{Not in the dynamical sense being used here.  Entanglement
creates correlations, but a local operation by Alice cannot be used to
alter Bob's reduced state or send him a controllable signal
instantaneously.  The nickel-and-dime analogy isolates the distinction:
seeing one coin reveals which correlated coin the other person has, but it
does not act on the distant coin.  Quantum correlations are richer than
that classical analogy, yet they still do not provide superluminal
signaling. \lecturetimestamp{00:20:02}}
\end{classroomqa}

\subsection{Lorentz invariance}

The third demand is that the laws have the same form in every inertial
frame.  The practical construction rule is simple: the Lagrangian must be
a Lorentz scalar, so the action is invariant.  Upper--lower contractions
such as \(\partial_\mu\phi\,\partial^\mu\phi\) provide the building blocks.

In general relativity the statement is enlarged.  The integrand is a scalar
density, and invariance extends from Lorentz transformations to arbitrary
coordinate transformations.  That refinement does not replace the present
principle; it is its curved-spacetime form.

\begin{classroomqa}
\audiencequestion{Is Lorentz invariance only a rule for special relativity, rather than a truly general principle?}
\lecturerresponse{For a nongravitational problem in flat spacetime, Lorentz
invariance is the correct statement.  In curved spacetime one replaces the
scalar Lagrangian by the appropriate scalar density and requires invariance
under arbitrary coordinate transformations.  Apart from that refinement,
the same construction philosophy survives. \lecturetimestamp{00:24:16}}
\end{classroomqa}

\begin{classroomqa}
\audiencequestion{What are the dimensions of the particle Lagrangian and the field Lagrangian density?}
\lecturerresponse{A particle Lagrangian has dimensions of energy.  In
\(I=\int dt\,d^3x\,\mathcal L\), the spatial measure supplies a volume, so
\(\mathcal L\) has dimensions of energy per unit volume.  Integrating it
over space gives an ordinary Lagrangian with dimensions of energy.
\lecturetimestamp{00:25:49}}
\end{classroomqa}

\subsection{Gauge invariance}

The fourth demand is gauge invariance.  At this point it is deliberately
left mysterious.  An audience member offers a useful anticipation: perhaps
one can add certain things to a Lagrangian without changing the physics.
That is close, but the word certain will do nearly all the work.

\begin{figure}[htbp]
\centering
\includegraphics[width=0.90\textwidth]{lecture_07_figure_04.png}
\caption{Locality, scalar construction, Lorentz invariance, and gauge invariance collected on the board.}
\lectureframe{00:27:00}
\end{figure}

We postpone the definition because the charged-particle calculation already
contains it.  Repeating that calculation is not a detour.  It will show why
the vector potential is forced upon us and why, once introduced, it contains
a redundancy.

\section{Why the Vector Potential Is Necessary}

For a slowly moving charged particle, the target equation is the familiar
Lorentz-force law
\begin{equation}
m\mathbf a
=
e\left(
\mathbf E+\mathbf v\times\mathbf B
\right).
\label{eq:sr7-lorentz-target}
\end{equation}

\begin{figure}[htbp]
\centering
\includegraphics[width=0.90\textwidth]{lecture_07_figure_05.png}
\caption{The familiar force law, written before the vector potential is reintroduced.}
\lectureframe{00:29:08.870}
\end{figure}

No vector potential appears in \eqref{eq:sr7-lorentz-target}.  One can write
the equations of motion directly with \(\mathbf E\) and \(\mathbf B\).
The difficulty is more structural: under the assumptions of locality and
Lorentz invariance used in the lecture, the interaction of a charged
particle cannot be written as a worldline action directly in terms of those
fields, position, and velocity alone.  The potential \(A_\mu\) is needed to
construct the local scalar coupling.

\begin{classroomqa}
\audiencequestion{If the force law already contains only \(\mathbf E\) and \(\mathbf B\), why introduce \(A_\mu\)?}
\lecturerresponse{Because an equation of motion is not yet an action
principle.  Within the local Lorentz-invariant framework of this lecture,
the charged-particle action requires the four-potential.  The fields alone
can state the force, but the potential is the object that permits the force
to be derived from a scalar worldline action. \lecturetimestamp{00:29:06}}
\end{classroomqa}

\subsection{The worldline action}

The candidate action has a free part and an interaction part:
\begin{equation}
I[x]
=
-m\int d\tau
-e\int A_\mu(x)\,dx^\mu.
\label{eq:sr7-worldline-action}
\end{equation}
The first term adds the proper-time length of every infinitesimal segment of
the trajectory.  It is a scalar by construction.  In the second term,
\(A_\mu(x)\) is evaluated at the event occupied by the particle and
contracted with the displacement \(dx^\mu\).  That increment is also a
scalar and is local along the worldline.

\begin{figure}[htbp]
\centering
\includegraphics[width=0.90\textwidth]{lecture_07_figure_06.png}
\caption{The two pieces of the charged-particle action and the local contraction \(A_\mu(x)dx^\mu\).}
\lectureframe{00:31:20}
\end{figure}

\subsection{Rewriting the action in ordinary time}

To use the familiar Euler--Lagrange equations, divide each worldline
increment by \(dt\):
\begin{equation}
\frac{d\tau}{dt}
=
\sqrt{1-\dot{\mathbf x}^{\,2}},
\qquad
\dot{\mathbf x}^{\,2}
=
\dot x^2+\dot y^2+\dot z^2.
\end{equation}
The four components of \(dx^\mu/dt\) are
\begin{equation}
\frac{dx^\mu}{dt}
=
\left(1,\dot x,\dot y,\dot z\right).
\end{equation}
Consequently,
\begin{align}
I
&=
\int dt\,
\left[
-m\sqrt{1-\dot{\mathbf x}^{\,2}}
-eA_\mu\frac{dx^\mu}{dt}
\right]
\notag\\
&=
\int dt\,
\left[
-m\sqrt{1-\dot{\mathbf x}^{\,2}}
-e\left(A_0+A_m\dot x^m\right)
\right].
\label{eq:sr7-action-in-t}
\end{align}
The ordinary-time Lagrangian is therefore
\begin{equation}
\begin{aligned}
L(x,\dot x,t)
&=
-m\sqrt{1-\dot{\mathbf x}^{\,2}}
\\
&\quad -e\left[
A_0(x,t)+A_m(x,t)\dot x^m
\right].
\end{aligned}
\label{eq:sr7-lagrangian}
\end{equation}

\begin{figure}[htbp]
\centering
\includegraphics[width=0.90\textwidth]{lecture_07_figure_07.png}
\caption{The worldline action rewritten as an ordinary-time integral and the Lagrangian read from its integrand.}
\lectureframe{00:35:30}
\end{figure}

At this stage the imaginative part is over.  Once a Lagrangian has been
chosen, the rest is mechanical.  The mechanics can be long and tedious,
but that tedium is not empty: it is what reveals the electric, magnetic, and
relativistic structures.  The lecturer jokes that discovering the principle
is the enjoyable part and deducing all its consequences is what one gives to
graduate students.

\begin{classroomqa}
\audiencequestion{May the spatial index on \(A_m\) or \(\dot x^m\) be moved upstairs or downstairs in this calculation?}
\lecturerresponse{Yes.  The spatial metric is the identity, so raising or
lowering a Latin index does not change its component value.  The time index
is the one whose sign requires care. \lecturetimestamp{00:34:10}}
\end{classroomqa}

\section{The Euler--Lagrange Calculation}

\subsection{Canonical momentum and a live correction}

Choose one spatial component \(x^m\).  Differentiating the free term of
\eqref{eq:sr7-lagrangian} gives
\begin{equation}
\frac{m\dot x_m}{\sqrt{1-\dot{\mathbf x}^{\,2}}}
=
m u_m,
\qquad
u^\mu\equiv\frac{dx^\mu}{d\tau}.
\end{equation}
The first board attempt stops there.  An audience correction catches the
omission: the interaction term also depends on \(\dot x^m\).  The complete
derivative is
\begin{equation}
\frac{\partial L}{\partial\dot x^m}
=
\frac{m\dot x_m}{\sqrt{1-\dot{\mathbf x}^{\,2}}}
-eA_m
=
m u_m-eA_m.
\label{eq:sr7-canonical-momentum}
\end{equation}

\begin{classroomqa}
\audiencequestion{Is the first expression for \(\partial L/\partial\dot x^m\) missing the contribution from the electromagnetic term?}
\lecturerresponse{Yes.  Differentiating
\(-eA_n\dot x^n\) with respect to \(\dot x^m\) contributes
\(-eA_m\).  The corrected canonical momentum is
\(m u_m-eA_m\), not the mechanical momentum alone.
\lecturetimestamp{00:37:24}}
\end{classroomqa}

\begin{figure}[htbp]
\centering
\includegraphics[width=0.90\textwidth]{lecture_07_figure_08.png}
\caption{The corrected canonical momentum with both the mechanical and potential terms.}
\lectureframe{00:38:30}
\end{figure}

\subsection{Differentiating along the moving path}

The left side of the Euler--Lagrange equation is
\begin{equation}
\frac{d}{dt}
\left(
\frac{\partial L}{\partial\dot x^m}
\right)
=
m\frac{du_m}{dt}
-e\frac{dA_m(x(t),t)}{dt}.
\end{equation}
The potential changes for two reasons: it may depend explicitly on time,
and the particle moves through its spatial variation.  The chain rule gives
\begin{equation}
\frac{dA_m}{dt}
=
\partial_0A_m
+\dot x^n\partial_nA_m.
\label{eq:sr7-chain-rule}
\end{equation}
In longhand, the second term means
\begin{equation}
\dot x\,\frac{\partial A_m}{\partial x}
+\dot y\,\frac{\partial A_m}{\partial y}
+\dot z\,\frac{\partial A_m}{\partial z}.
\end{equation}

\begin{figure}[htbp]
\centering
\includegraphics[width=0.90\textwidth]{lecture_07_figure_09.png}
\caption{The explicit and convective pieces of the time derivative of \(A_m\).}
\lectureframe{00:42:10}
\end{figure}

On the right side,
\begin{equation}
\frac{\partial L}{\partial x^m}
=
-e\partial_mA_0
-e\dot x^n\partial_mA_n.
\label{eq:sr7-spatial-derivative-L}
\end{equation}
The free term has no explicit position dependence.  In
\eqref{eq:sr7-spatial-derivative-L}, \(m\) is the free index selecting one
equation; \(n\) is a repeated dummy index and is summed.  Renaming that
dummy index changes nothing.  This distinction is why the lecturer pauses
over the letters before combining the terms.

Equating the two sides and rearranging yields
\begin{equation}
m\frac{du_m}{dt}
=
e\left(
\partial_0A_m-\partial_mA_0
\right)
+e\dot x^n
\left(
\partial_nA_m-\partial_mA_n
\right).
\label{eq:sr7-grouped-force}
\end{equation}

\begin{figure}[htbp]
\centering
\includegraphics[width=0.90\textwidth]{lecture_07_figure_10.png}
\caption{The force separated into a velocity-independent bracket and a velocity-linear antisymmetric bracket.}
\lectureframe{00:49:30}
\end{figure}

\subsection{Recognizing electric and magnetic fields}

The first bracket contains no particle velocity.  In the convention of this
action, define
\begin{equation}
E_m
\equiv
\partial_0A_m-\partial_mA_0.
\label{eq:sr7-E-from-A}
\end{equation}
It has precisely the operational role of the electric field.  The second
bracket is antisymmetric in its spatial indices:
\begin{equation}
C_{mn}
\equiv
\partial_nA_m-\partial_mA_n.
\label{eq:sr7-C}
\end{equation}
Its three independent components are the curl of the spatial potential,
with orientation chosen in the lecture convention so that
\begin{equation}
\dot x^n C_{mn}
=
\left(\mathbf v\times\mathbf B\right)_m.
\label{eq:sr7-magnetic-identification}
\end{equation}
The detailed component check is left as an exercise on the board.  The
structural identification is already visible: an antisymmetric spatial
derivative of a vector has three components, exactly the number carried by
\(\mathbf B\), and it multiplies the velocity.

\begin{figure}[htbp]
\centering
\includegraphics[width=0.90\textwidth]{lecture_07_figure_11.png}
\caption{The two derivative brackets identified as the electric and magnetic parts of the Lorentz force.}
\lectureframe{00:52:00}
\end{figure}

Equation \eqref{eq:sr7-grouped-force} has therefore recovered the Lorentz
force from the vector-potential action.  It has also done something more
important: it has exposed the antisymmetrized derivatives that will become
the field tensor.

\section{From Ordinary Time to a Four-Vector Equation}

\subsection{The proper-time conversion}

The left side of \eqref{eq:sr7-grouped-force} is not the Newtonian
acceleration \(d\mathbf v/dt\).  It is the ordinary-time derivative of a
four-velocity component:
\begin{equation}
\frac{du^m}{dt}
=
\frac{du^m}{d\tau}\frac{d\tau}{dt}.
\end{equation}
Multiply the equation by \(dt/d\tau\).  Then
\begin{equation}
\dot x^n\frac{dt}{d\tau}
=
\frac{dx^n}{dt}\frac{dt}{d\tau}
=
u^n,
\qquad
\frac{dt}{d\tau}=u^0.
\end{equation}
The spatial equation becomes
\begin{align}
m\frac{du^m}{d\tau}
&=
e\left(
\partial_0A_m-\partial_mA_0
\right)u^0
\notag\\
&\quad
+e\left(
\partial_nA_m-\partial_mA_n
\right)u^n.
\label{eq:sr7-proper-spatial}
\end{align}

\begin{figure}[htbp]
\centering
\includegraphics[width=0.90\textwidth]{lecture_07_figure_12.png}
\caption{The factors \(dt/d\tau\) convert ordinary velocities into components of the four-velocity.}
\lectureframe{00:56:30}
\end{figure}

The left side is proper acceleration:
\begin{equation}
\frac{du^\mu}{d\tau}
=
\frac{d^2x^\mu}{d\tau^2}.
\end{equation}
It is a four-vector.  The right side must therefore be assembled from a
rank-two tensor acting on \(u^\nu\).

\subsection{The field tensor and the force law}

Choose the convention
\begin{equation}
F_{\mu\nu}
\equiv
\partial_\mu A_\nu-\partial_\nu A_\mu.
\label{eq:sr7-F-definition}
\end{equation}
It is antisymmetric:
\begin{equation}
F_{\mu\nu}=-F_{\nu\mu}.
\end{equation}
Using the interaction sign in \eqref{eq:sr7-worldline-action}, the compact
equation is
\begin{equation}
\boxed{
m\frac{du^\mu}{d\tau}
=
-eF^\mu{}_\nu u^\nu
}.
\label{eq:sr7-covariant-force}
\end{equation}
Equivalently, before naming \(F_{\mu\nu}\),
\begin{equation}
m\frac{du^\mu}{d\tau}
=
e\left(
\partial_\nu A^\mu-\partial^\mu A_\nu
\right)u^\nu.
\end{equation}
Different metric, charge, and potential conventions redistribute the
overall signs.  The invariant content is the antisymmetrized derivative of
the potential and its contraction with the four-velocity.

\begin{figure}[htbp]
\centering
\includegraphics[width=0.90\textwidth]{lecture_07_figure_13.png}
\caption{The proper-time equation packaged as a rank-two field tensor acting on the four-velocity.}
\lectureframe{00:58:25}
\end{figure}

\begin{classroomqa}
\audiencequestion{The derivation produced three spatial equations.  Why is a fourth equation allowed to be added?}
\lecturerresponse{The scalar action guarantees a Lorentz-covariant result.
Once the difference between the two sides is recognized as a four-vector,
Lorentz transformations mix its spatial and temporal components.  A
four-vector equation cannot consistently have only three components; the
time component accompanies the spatial ones. \lecturetimestamp{00:59:40}}
\end{classroomqa}

\begin{classroomqa}
\audiencequestion{Is a fixed list of numbers such as \((1,2,3,4)\) itself a Lorentz-invariant four-vector?}
\lecturerresponse{No.  Such numbers can be the components of a vector in
one chosen frame, but the components transform.  Even in two dimensions,
the vector with components \((1,2)\) can be described in a rotated frame as
\((0,\sqrt5)\).  The geometric vector persists; its numerical component
list does not. \lecturetimestamp{01:00:28}}
\end{classroomqa}

\begin{figure}[htbp]
\centering
\includegraphics[width=0.90\textwidth]{lecture_07_figure_14.png}
\caption{A two-component vector re-expressed in axes aligned with the vector itself.}
\lectureframe{01:02:00}
\end{figure}

\begin{classroomqa}
\audiencequestion{If the goal was to replace ordinary time by proper time, why does \(\partial_0A_m=\partial A_m/\partial t\) still appear?}
\lecturerresponse{The goal was not to erase every letter \(t\).  It was to
write the dynamics in manifest four-vector notation.  The operator
\(\partial_0\) is the time component of the covariant four-gradient, just as
\(\partial_m\) are its spatial components.  It belongs inside the
four-dimensional tensor. \lecturetimestamp{01:05:10}}
\end{classroomqa}

The lecture uses this moment to restate locality.  A local equation relates
the acceleration at an event to the field and velocity at that event.  It
does not link the acceleration now to a force at a remote time.

\begin{figure}[htbp]
\centering
\includegraphics[width=0.90\textwidth]{lecture_07_figure_15.png}
\caption{The covariant force equation used to restate locality as differential dynamics.}
\lectureframe{01:05:50}
\end{figure}

\subsection{Displaying the electromagnetic components}

For the all-upper tensor, the convention used on the board gives
\begin{equation}
F^{\mu\nu}
=
\begin{pmatrix}
0&-E_x&-E_y&-E_z\\
E_x&0&B_z&-B_y\\
E_y&-B_z&0&B_x\\
E_z&B_y&-B_x&0
\end{pmatrix}.
\label{eq:sr7-F-matrix}
\end{equation}
The row index is \(\mu\), the column index is \(\nu\).  The mixed
time--space entries contain the electric field; the space--space block
contains the magnetic field.  Antisymmetry sets the diagonal to zero and
reduces sixteen displayed entries to six independent components.

\begin{figure}[htbp]
\centering
\includegraphics[width=0.90\textwidth]{lecture_07_figure_16.png}
\caption{The completed electromagnetic tensor, with electric entries in the time--space slots and magnetic entries in the spatial block.}
\lectureframe{01:12:10}
\end{figure}

\begin{classroomqa}
\audiencequestion{Is there a standard rule deciding which tensor index labels rows and which labels columns?}
\lecturerresponse{It is a display convention.  Here the first index
\(\mu\) runs vertically and the second index \(\nu\) runs horizontally.
Once that convention is announced, every component can be read
unambiguously.  Raising a time index introduces a sign, but it does not
alter the row--column convention. \lecturetimestamp{01:09:01}}
\end{classroomqa}

\begin{classroomqa}
\audiencequestion{Which index placements preserve antisymmetry?}
\lecturerresponse{The tensors \(F_{\mu\nu}\) and \(F^{\mu\nu}\) are
antisymmetric.  The mixed array \(F^\mu{}_\nu\) is not antisymmetric under
ordinary matrix transposition: raising one time index changes the electric
slots, making the time--space pair symmetric, while the purely spatial
magnetic block remains antisymmetric.  The lecturer reaches this conclusion
after briefly testing and rejecting the opposite possibility on the board.
\lecturetimestamp{01:13:18}}
\end{classroomqa}

\section{What Counts as an Invariance?}

Only now do we return to the fourth principle.  An invariance is a change
that leaves the equations of motion unchanged.  The action itself may remain
exactly unchanged, or it may change by a boundary term that has no effect on
the variational equations.

\subsection{Constant shifts of a massless scalar}

Consider the massless scalar equation in one spatial dimension,
\begin{equation}
\frac{\partial^2\phi}{\partial t^2}
-\frac{\partial^2\phi}{\partial x^2}
=0,
\label{eq:sr7-wave}
\end{equation}
obtained from
\begin{equation}
\mathcal L_0
=
-\frac12\partial_\mu\phi\,\partial^\mu\phi.
\label{eq:sr7-massless-L}
\end{equation}
If \(\phi\) is a solution, then so is
\begin{equation}
\phi\longrightarrow\phi+c,
\label{eq:sr7-constant-shift}
\end{equation}
where \(c\) is constant.  Every derivative of \(c\) vanishes, so neither the
equation nor the derivative-only Lagrangian changes.

\begin{figure}[htbp]
\centering
\includegraphics[width=0.90\textwidth]{lecture_07_figure_17.png}
\caption{The massless wave equation, its derivative Lagrangian, and the constant shift \(\phi\mapsto\phi+c\).}
\lectureframe{01:17:10}
\end{figure}

Now add a mass term:
\begin{equation}
\frac{\partial^2\phi}{\partial t^2}
-\frac{\partial^2\phi}{\partial x^2}
+\mu^2\phi
=0,
\end{equation}
\begin{equation}
\mathcal L_\mu
=
-\frac12\partial_\mu\phi\,\partial^\mu\phi
-\frac{\mu^2}{2}\phi^2.
\label{eq:sr7-massive-L}
\end{equation}
The derivative part remains unchanged under
\(\phi\mapsto\phi+c\), but \(\phi^2\) does not.  Both theories are
legitimate; one has the shift invariance and the other does not.

\begin{figure}[htbp]
\centering
\includegraphics[width=0.90\textwidth]{lecture_07_figure_18.png}
\caption{The mass term that breaks constant-shift invariance.}
\lectureframe{01:18:50}
\end{figure}

\subsection{A student anticipates the boundary term}

An audience member proposes a sharper test:
\begin{equation}
\phi(x,t)\longrightarrow\phi(x,t)+x.
\label{eq:sr7-linear-shift}
\end{equation}
The second derivatives of \(x\) vanish, so the massless wave equation is
still satisfied.  The action requires a closer look.  Since \(x\) and \(t\)
are independent variables,
\begin{equation}
\frac{\partial x}{\partial t}=0,
\qquad
\frac{\partial x}{\partial x}=1.
\end{equation}
The spatial part changes according to
\begin{align}
-\frac12
\left[
\partial_x(\phi+x)
\right]^2
&=
-\frac12
\left(
\partial_x\phi+1
\right)^2
\notag\\
&=
-\frac12(\partial_x\phi)^2
-\partial_x\phi
-\frac12.
\label{eq:sr7-linear-shift-L}
\end{align}
The new pieces are a constant and a total derivative.  The constant does
not change the Euler--Lagrange equation.  Integrating
\(\partial_x\phi\) produces only spatial boundary values, so under the
usual fixed-boundary conditions it also leaves the equation of motion
unchanged.

\begin{figure}[htbp]
\centering
\includegraphics[width=0.90\textwidth]{lecture_07_figure_19.png}
\caption{The proposed shift \(\phi\mapsto\phi+x\) and the derivative square that reveals a constant plus a total derivative.}
\lectureframe{01:22:40}
\end{figure}

\begin{classroomqa}
\audiencequestion{Does \(\phi\mapsto\phi+x\) preserve the massless field equation even though it changes the Lagrangian density?}
\lecturerresponse{Yes.  The second derivatives of \(x\) vanish, so the
equation is unchanged.  The Lagrangian changes only by a constant and a
total spatial derivative.  These alter the numerical action by a fixed
bulk constant and boundary values, but not the Euler--Lagrange equation
under fixed boundary conditions. \lecturetimestamp{01:20:00}}
\end{classroomqa}

\begin{classroomqa}
\audiencequestion{Why is \(\partial x/\partial t\) zero?}
\lecturerresponse{Because \(x\) and \(t\) are independent coordinates of
the field description.  A partial derivative with respect to \(t\) holds
\(x\) fixed. \lecturetimestamp{01:22:22}}
\end{classroomqa}

The student has arrived at the device needed for electromagnetism: a total
derivative can change the written action without changing its stationary
paths.

\section{Gauge Invariance of the Charged Particle}

\subsection{Adding a four-gradient}

Return to the interaction
\begin{equation}
I_{\mathrm{int}}
=
-e\int A_\mu\,dx^\mu.
\end{equation}
Let \(\Lambda(x)\) be any sufficiently smooth scalar and change the
potential by
\begin{equation}
A_\mu
\longrightarrow
A'_\mu
=
A_\mu+\partial_\mu\Lambda.
\label{eq:sr7-gauge-transformation}
\end{equation}
The gradient \(\partial_\mu\Lambda\) has the same covariant index structure
as \(A_\mu\), so the sum is well defined.

\begin{figure}[htbp]
\centering
\includegraphics[width=0.90\textwidth]{lecture_07_figure_20.png}
\caption{The worldline coupling and the transformation \(A_\mu\mapsto A_\mu+\partial_\mu\Lambda\).}
\lectureframe{01:24:50}
\end{figure}

The change in the action is
\begin{align}
\Delta I_{\mathrm{int}}
&=
-e\int \partial_\mu\Lambda\,dx^\mu
\notag\\
&=
-e\int d\Lambda
\notag\\
&=
-e\left[
\Lambda(x_2)-\Lambda(x_1)
\right].
\label{eq:sr7-gauge-boundary}
\end{align}
Along each infinitesimal worldline segment,
\(\partial_\mu\Lambda\,dx^\mu\) is precisely the increment
\(d\Lambda\).  Adding the increments telescopes to the difference between
the final and initial endpoint values.

\begin{figure}[htbp]
\centering
\includegraphics[width=0.90\textwidth]{lecture_07_figure_21.png}
\caption{The gradient contribution integrated along the path to the endpoint difference \(\Lambda(x_2)-\Lambda(x_1)\).}
\lectureframe{01:28:10}
\end{figure}

The variational principle compares nearby trajectories with \(x_1\) and
\(x_2\) held fixed.  The extra term in \eqref{eq:sr7-gauge-boundary} is
therefore the same for every allowed variation:
\begin{equation}
\delta\Delta I_{\mathrm{int}}=0.
\end{equation}
The stationary trajectory is unchanged.  This is gauge invariance as an
endpoint statement.

\begin{classroomqa}
\audiencequestion{May \(\Lambda\) be any scalar function?}
\lecturerresponse{Yes, subject to the smoothness and boundary assumptions
needed for the derivatives and action.  The transformation is local because
the parameter is an arbitrary function of spacetime, not merely one global
constant. \lecturetimestamp{01:28:40}}
\end{classroomqa}

Adding a gradient to \(A_\mu\) is called a gauge transformation.  The word
gauge is a historical inheritance; the lecture warns, humorously but
firmly, against pronouncing it as though it were a French word.  The name is
less important than the operation.

\subsection{The direct proof with the field tensor}

The equations of motion depend on \(F_{\mu\nu}\), not on the potential by
itself.  Under \eqref{eq:sr7-gauge-transformation},
\begin{align}
F'_{\mu\nu}
&=
\partial_\mu
\left(
A_\nu+\partial_\nu\Lambda
\right)
-
\partial_\nu
\left(
A_\mu+\partial_\mu\Lambda
\right)
\notag\\
&=
F_{\mu\nu}
+
\partial_\mu\partial_\nu\Lambda
-
\partial_\nu\partial_\mu\Lambda
\notag\\
&=
F_{\mu\nu}.
\label{eq:sr7-F-gauge-invariant}
\end{align}
For a smooth scalar in ordinary flat spacetime, partial derivatives commute.
The two additional terms cancel.  Thus the field tensor, the electric and
magnetic fields it contains, and the force law are all gauge invariant.

\begin{figure}[htbp]
\centering
\includegraphics[width=0.90\textwidth]{lecture_07_figure_22.png}
\caption{The two second derivatives generated by a gauge transformation cancel because partial derivatives commute.}
\lectureframe{01:33:50}
\end{figure}

The two proofs close the same circle from opposite directions.  The action
changes only at fixed endpoints, while the equations depend only on an
unchanged field tensor.

\section{What the Redundancy Means}

\subsection{Why keep an object that is not unique?}

The natural objection is immediate.  If \(A_\mu\) is not unique, why did we
introduce it?  Because it is indispensable to the local Lorentz-invariant
action even though its component values contain redundant information.
Many potentials represent the same electromagnetic field:
\begin{equation}
A_\mu
\sim
A_\mu+\partial_\mu\Lambda.
\end{equation}
The equivalence class is physical; a particular representative is a choice
of description.

\begin{classroomqa}
\audiencequestion{Is there practical value in changing the gauge?}
\lecturerresponse{Yes.  A suitable \(\Lambda\) can simplify a component of
the potential or make one feature of a problem transparent.  Locally, and
with suitable boundary conditions, one can often choose a gauge in which
one selected component, including \(A_0\), vanishes.  Another gauge may
display a different feature more clearly. \lecturetimestamp{01:35:43}}
\end{classroomqa}

\begin{classroomqa}
\audiencequestion{Is gauge invariance analogous to describing the same experiment from two Lorentz frames?}
\lecturerresponse{No.  Lorentz transformations compare physically
meaningful reference frames carried by different observers.  Gauge
transformations are redundancies within one description: different
potentials encode the same field.  Mathematically both are invariances, but
their interpretations differ. \lecturetimestamp{01:37:07}}
\end{classroomqa}

There is also a structural difference.  A rotation uses a fixed global
angle; one does not ordinarily rotate the axes by unrelated angles at each
point.  A gauge transformation is parameterized by the entire function
\(\Lambda(x)\).  This position-dependent freedom is what makes it a local
gauge invariance.

\begin{classroomqa}
\audiencequestion{Is adding a gradient the general definition of gauge invariance?}
\lecturerresponse{For the Abelian electromagnetic theory, that is the
transformation.  Yang--Mills theories have a more elaborate non-Abelian
version in which the potential carries additional internal structure, but
the same idea survives: different local descriptions represent the same
physics. \lecturetimestamp{01:39:27}}
\end{classroomqa}

\begin{classroomqa}
\audiencequestion{Is a gauge transformation merely a mathematical device for highlighting different parts of a problem?}
\lecturerresponse{It is a useful calculational and interpretive device, but
its role is constrained: it may simplify the equations or reveal a
property, yet it cannot change any final gauge-invariant prediction.
\lecturetimestamp{01:40:22}}
\end{classroomqa}

\subsection{Other fields and other potentials}

Gauge potentials in other theories are also commonly denoted by \(A\) and
carry a spacetime index \(\mu\).  What distinguishes them is additional
internal structure associated with other symmetries.  The lecture postpones
that structure rather than offering a premature formula.

\begin{classroomqa}
\audiencequestion{What do the vector potentials of other field theories look like?}
\lecturerresponse{They still behave as spacetime vectors and are commonly
written \(A_\mu\), but they carry additional labels and algebraic structure
associated with internal symmetries.  That is the entrance to Yang--Mills
theory and is intentionally deferred. \lecturetimestamp{01:41:02}}
\end{classroomqa}

\begin{classroomqa}
\audiencequestion{Is it really impossible to write the local Lorentz-invariant charged-particle action directly with \(\mathbf E\) and \(\mathbf B\)?}
\lecturerresponse{Within the assumptions stated here, yes.  The lecture
presents this as a theorem and does not prove it.  Preserving the action
principle, locality, and Lorentz invariance forces the introduction of the
potential, after which gauge redundancy appears automatically.
\lecturetimestamp{01:42:18}}
\end{classroomqa}

\subsection{What can be measured?}

A local value of \(A_\mu\) by itself is not gauge invariant.  If a physical
measurement determined that value uniquely, changing it by
\(\partial_\mu\Lambda\) would alter a physical outcome, contrary to the
invariance just proved.  Measurable predictions must be built from
gauge-invariant quantities, such as \(F_{\mu\nu}\) in the present classical
discussion.

\begin{classroomqa}
\audiencequestion{Why does nature so persistently choose gauge theories?}
\lecturerresponse{At the classical level the lecture offers no final reason.
Quantum mechanics supplies a deeper perspective, but even there the broad
answer is not exhausted.  What is already evident is the unusual beauty,
symmetry, and economy of theories organized by gauge redundancy.
\lecturetimestamp{01:43:36}}
\end{classroomqa}

\begin{classroomqa}
\audiencequestion{Can the vector potential itself be measured?}
\lecturerresponse{Not as a gauge-dependent local component value.  A
measurement is a physical process, and every physical prediction must be
insensitive to \(A_\mu\mapsto A_\mu+\partial_\mu\Lambda\).  What can be
measured are gauge-invariant features of the potential and field.
\lecturetimestamp{01:44:18}}
\end{classroomqa}

\begin{classroomqa}
\audiencequestion{What happens if \(A_\mu\) is changed by something that is not a gradient?}
\lecturerresponse{In general \(F_{\mu\nu}\) changes, so the change is
detectable.  By contrast, a pure-gauge potential
\(A_\mu=\partial_\mu\Lambda\) has
\(F_{\mu\nu}=0\) and exerts no electromagnetic force; a free particle then
moves on a straight worldline. \lecturetimestamp{01:45:22}}
\end{classroomqa}

Gauge invariance is therefore not a disposable trick.  Electrodynamics,
Yang--Mills theory, the Standard Model, and gravity all possess gauge
redundancies of one kind or another.  The pattern is peculiar: an auxiliary
quantity is required to formulate the theory, yet part of that quantity can
be changed without changing the physics.  The peculiarity is also the
organizing power.

\section{The Logic of Gauge Invariance}

The lecture has taken a deliberately indirect path.  We first restored the
upper--lower index language because scalar contractions are the grammar of
Lorentz-invariant actions.  We then stated four principles: action,
locality, Lorentz invariance, and gauge invariance.  The fourth remained
undefined while the charged-particle calculation was rebuilt.

That calculation began with the local scalar action
\begin{equation}
I
=
-m\int d\tau
-e\int A_\mu dx^\mu,
\end{equation}
passed through the ordinary-time Lagrangian and every Euler--Lagrange
derivative, and recovered the electric and magnetic Lorentz force.  Proper
time then assembled the result into
\begin{equation}
m\frac{du^\mu}{d\tau}
=
-eF^\mu{}_\nu u^\nu,
\qquad
F_{\mu\nu}
=
\partial_\mu A_\nu-\partial_\nu A_\mu.
\end{equation}

The scalar-field examples prepared the final step: a total derivative can
alter an action only through boundaries.  Consequently,
\begin{equation}
A_\mu\longrightarrow A_\mu+\partial_\mu\Lambda
\end{equation}
changes the worldline action only by fixed endpoint values and leaves
\(F_{\mu\nu}\) unchanged because mixed partial derivatives commute.  The
vector potential is therefore necessary but nonunique.  Gauge invariance is
the statement that this nonuniqueness belongs to the description, not to
the physics.
